\chapter{Related Work}
\label{chap:related}

Existing work reduces what the LLM reads and retrieves evidence iteratively, but it leaves open whether a very small selector can learn to use the evidence it has already selected, and how that ability is represented in its weights. This chapter reviews three lines of research: context selection for RAG, multi-hop retrieval, and the structure of fine-tuning updates. Table~\ref{tab:positioning} summarizes what is known and what this work adds.

\section{Context selection and compression for RAG}

Several methods reduce what the LLM reads by compressing the context. RECOMP \cite{xu2024recomp} trains compressors that summarize retrieved documents; LongLLMLingua \cite{jiang2024longllmlingua} prunes tokens; FILCO \cite{wang2023filco} filters context at the sentence level; Provence \cite{chirkova2025provence} prunes sentences with a trained cross-encoder; and CompAct \cite{yoon2024compact} compresses iteratively.

Other methods decide how much to retrieve or send. Self-RAG \cite{asai2024selfrag} and Adaptive-RAG \cite{jeong2024adaptiverag} use the LLM or a classifier to decide, DynamicRAG \cite{sun2025dynamicrag} learns the number of documents, and Adaptive-$k$ \cite{taguchi2025adaptivek} cuts the ranked list at the largest gap in similarity scores and adds a buffer of extra passages (five by default). Noise-Gate \cite{lavado2027noisegate}, the author's earlier work, selects paragraphs by embedding-similarity thresholds and reported 26--27\% fewer tokens on HotpotQA; in that study Adaptive-$k$ selected too few passages and lowered answer quality. Listwise LLM rerankers \cite{sun2023rankgpt,reddy2024first} are accurate but require an LLM call per ranked list, far more computation than a small cross-encoder.

The present work does not propose a new compression mechanism. Its contribution to this line is a cost accounting of a very small selector at fixed recall targets, with answer-F1 differences and confidence intervals, and the observation that on harder multi-hop data the quality of the ranking, rather than the stopping rule, limits cost.

\section{Multi-hop retrieval}

Multi-hop dense retrieval \cite{xiong2021mdr} and Baleen \cite{khattab2021baleen} retrieve iteratively, conditioning each step on what was found before. IRCoT \cite{trivedi2023ircot} interleaves retrieval with chain-of-thought reasoning by an LLM, and EfficientRAG \cite{zhuang2024efficientrag} replaces some LLM calls with small models. Beam Retrieval \cite{zhang2024beamretrieval} is the method closest to the iterative selector studied here: a cross-encoder scores each candidate together with the question and the passages already chosen, with beam search. It reports high supporting-passage F1 on MuSiQue with much larger encoders (DeBERTa base and large).

The iterative selector of this work follows this design. The contributions here are (i) its evaluation at 22.7M parameters with a pair-matched control that separates the effect of conditioning from the effect of more training data, and (ii) the analysis of the weight change that conditioning requires.

\section{Structure of fine-tuning updates}

The low intrinsic dimension of fine-tuning \cite{li2018intrinsic,aghajanyan2021intrinsic} and the success of low-rank adapters \cite{hu2022lora} are well established. PHLoRA \cite{vasani2025phlora} extracts low-rank adapters from full fine-tuning checkpoints with truncated SVD, the same operation used here. Skill-localization studies find that a very small fraction of the parameters can carry a fine-tuned skill \cite{panigrahi2023skilllocalization}. For rerankers, analyses of LLMs fine-tuned with LoRA locate much of the relevance gain in the attention of middle layers \cite{perlman2026emergence}, and report that LoRA ranks as low as one already give most of the reranking quality \cite{nijasure2025howrelevance}; in a MiniLM cross-encoder, specific attention heads compute a BM25-like relevance signal \cite{lu2025pathway}. The change studied in this work is distributed more evenly over the layers than the middle-layer locus reported for LLM rerankers (Section~\ref{sec:res_rank}). LoRA updates differ structurally from full fine-tuning \cite{biderman2024loralearnsless,shuttleworth2025illusion}, and rank measures of updates need a noise baseline before they are interpreted \cite{altaf2026noisebaseline}.

Behaviours, not only domains, can be carried by very low-rank updates. A single rank-1 adapter can induce emergent misalignment in small LLMs \cite{turner2025modelorganisms}, and in a 14B-parameter model the misalignment induced by rank-1 adapters is mediated by an activation direction whose ablation also removes it in other fine-tunes \cite{soligo2025convergent}. Reasoning gains from reinforcement learning can be recovered with tiny adapters \cite{morris2026tinylora,ye2026rlvrlowrank}, while supervised fine-tuning needs much larger updates \cite{morris2026tinylora}. Task arithmetic merges domain and task weights to adapt LLM rerankers to new domains without further training \cite{braga2025taskarithmetic}.

Prior work therefore leads one to expect low-rank fine-tuning updates, including low-rank adapters for reranking \cite{nijasure2025howrelevance}. Those adapters are trained for relevance itself, whereas the change studied here is extracted from a full fine-tune and isolates one added behaviour. The claim tested here is more specific: a \emph{retrieval behaviour}, the use of already-selected evidence in the query, is carried in a small selector by a change of functional rank at most two. This is shown with a pair-matched control, norm-matched random controls and a seed-noise control, on three datasets. The present evidence concerns the sufficiency of a weight-space truncation; whether the change is mediated by a specific activation direction, and whether it transfers between models, is left to the remaining work (Chapter~\ref{chap:plan}).

\begin{table}[htbp]
\centering
\caption{Positioning: what is known and what this work adds.}
\label{tab:positioning}
\small
\begin{tabularx}{\textwidth}{>{\raggedright\arraybackslash}p{3.4cm}X>{\raggedright\arraybackslash}p{4.4cm}}
\toprule
Topic & Known & Added here \\
\midrule
Small selector reduces tokens & Trained selectors and compressors reduce tokens with little or no quality loss \cite{xu2024recomp,chirkova2025provence} & Token accounting at fixed recall targets for a 22.7M-parameter selector \\
Adaptive stopping & Adaptive-$k$, DynamicRAG, Self-RAG \cite{taguchi2025adaptivek,sun2025dynamicrag,asai2024selfrag} & Learned stopping on cross-encoder scores; Adaptive-$k$ reaches lower recall on these scores \\
Iterative selection & Beam Retrieval, MDR, Baleen \cite{zhang2024beamretrieval,xiong2021mdr,khattab2021baleen} & Pair-matched control: part of the gain came from more training data; the remaining gain concentrates on long chains \\
Low-rank updates & LoRA, PHLoRA, intrinsic dimension; rank-1 LoRA for LLM rerankers \cite{hu2022lora,vasani2025phlora,aghajanyan2021intrinsic,nijasure2025howrelevance} & A retrieval behaviour has functional rank at most two, with random, seed-noise and matched controls \\
Behaviour in tiny adapters & Rank-1 behaviours in LLMs, mediated by an activation direction \cite{soligo2025convergent} & A retrieval behaviour in a 22.7M-parameter cross-encoder restored by a rank-2 weight graft, on three datasets; mediation and transfer not yet tested \\
\bottomrule
\end{tabularx}
\end{table}
